% ============================================================
% Model: 自旋依赖势与低能散射
% ============================================================

\section{自旋依赖势与低能散射}
\label{sec:spin-dependent-potential}

\begin{tipbox}[关键词标签]
自旋依赖势 · 低能散射 · 全同费米子 · 分波法 · 自旋单态 · 自旋三重态 · 有效势
\end{tipbox}

% --------------------------------------------------------
% 1. 模型定义框
% --------------------------------------------------------
\begin{modelbox}[模型：低能中子-中子自旋依赖散射]
\textbf{物理情境}：两个中子发生低能散射，相互作用依赖于它们的相对自旋取向。由于中子是自旋-$1/2$ 费米子，总波函数必须反对称，这导致空间波函数和自旋波函数之间存在强烈的关联。

\textbf{相互作用势}：
\[
V(r)=\begin{cases}V_0\,\boldsymbol{\sigma}_1\cdot\boldsymbol{\sigma}_2,&r<a\\0,&r\geq a\end{cases}
\]

\textbf{自旋耦合分解}：
\begin{itemize}[nosep]
    \item 总自旋 $\mathbf{S}=\mathbf{S}_1+\mathbf{S}_2=\frac{\hbar}{2}(\boldsymbol{\sigma}_1+\boldsymbol{\sigma}_2)$
    \item $\boldsymbol{\sigma}_1\cdot\boldsymbol{\sigma}_2=\frac{2\mathbf{S}^2}{\hbar^2}-3$
    \item 单态（$S=0$）：$\boldsymbol{\sigma}_1\cdot\boldsymbol{\sigma}_2=-3$
    \item 三重态（$S=1$）：$\boldsymbol{\sigma}_1\cdot\boldsymbol{\sigma}_2=+1$
\end{itemize}

\textbf{有效势}：
\[
V_{\text{eff}}(r)=\begin{cases}-3V_0&(S=0,\text{ 单态})\\+V_0&(S=1,\text{ 三重态})\end{cases}
\]
\end{modelbox}

% --------------------------------------------------------
% 2. 关键公式框
% --------------------------------------------------------
\begin{formulabox}[核心公式]
\begin{enumerate}[label=\textbf{(\arabic*)}]
    \item \textbf{自旋态构造}：
    \[
    |\text{单态}\rangle=\frac{1}{\sqrt2}(|\uparrow\downarrow\rangle-|\downarrow\uparrow\rangle),\quad S=0,\;S_z=0,
    \]
    \[
    |\text{三重态}\rangle=\begin{cases}|\uparrow\uparrow\rangle,&S_z=1\\\frac{1}{\sqrt2}(|\uparrow\downarrow\rangle+|\downarrow\uparrow\rangle),&S_z=0\\|\downarrow\downarrow\rangle,&S_z=-1\end{cases}\quad S=1.
    \]

    \item \textbf{非极化入射的统计权重}：
    \[
    P(S=0)=\frac14,\quad P(S=1)=\frac34.
    \]
    \texteq{单态 1 个，三重态 3 个，共 4 个自旋态，等概率}

    \item \textbf{全同费米子的散射振幅}：
    \[
    \sigma(\theta)=|f(\theta)\pm f(\pi-\theta)|^2,
    \]
    $+$ 对应玻色子，$-$ 对应费米子。对非极化中子-中子散射，需按自旋通道分别计算后加权平均。

    \item \textbf{s 波散射长度}（低能极限，$ka\ll1$）：
    \[
    f_{S=0}=-a_s,\quad f_{S=1}=-a_t,
    \]
    $a_s$ 和 $a_t$ 分别为单态和三重态散射长度，由势阱深度和宽度决定。

    \item \textbf{总截面}（非极化，低能）：
    \[
    \sigma_{\text{tot}}=\frac14\cdot4\pi a_s^2+\frac34\cdot4\pi a_t^2=\pi(a_s^2+3a_t^2).
    \]
    \texteq{s 波全同费米子：$\sigma(\theta)=|f(\theta)-f(\pi-\theta)|^2$，对 s 波 $f$ 为常数，$\sigma=0$（向前=向后），总截面由自旋单态贡献}

    \item \textbf{磁偶极相互作用}（双中子）：
    \[
    H_{\text{dipole}}=\frac{\boldsymbol{\mu}_1\cdot\boldsymbol{\mu}_2}{r^3}-3\frac{(\boldsymbol{\mu}_1\cdot\mathbf{r})(\boldsymbol{\mu}_2\cdot\mathbf{r})}{r^5},
    \]
    $\boldsymbol{\mu}=\mu\boldsymbol{\sigma}$，取向沿 $z$ 轴时化简为 $H\propto\boldsymbol{\sigma}_1\cdot\boldsymbol{\sigma}_2-3\sigma_{1z}\sigma_{2z}$。
\end{enumerate}
\end{formulabox}

% --------------------------------------------------------
% 3. 训练点框
% --------------------------------------------------------
\begin{drillbox}[训练点与考法变体]
\trainingpoint{1} \textbf{自旋通道分离}：给定自旋依赖势 $V(r)\,\boldsymbol{\sigma}_1\cdot\boldsymbol{\sigma}_2$，写出单态和三重态的有效势，分别求解散射问题。

\trainingpoint{2} \textbf{全同粒子统计的应用}：两个中子（费米子）的 s 波空间波函数对称，要求自旋必须单态（反对称）。三重态的 s 波散射被统计禁止。

\trainingpoint{3} \textbf{非极化平均}：入射束非极化时，单态概率 $1/4$，三重态概率 $3/4$。计算总截面时需按此权重平均各通道的截面。

\trainingpoint{4} \textbf{低能极限与散射长度}：$k\to0$ 时仅 s 波有贡献，$f\to-a$。散射长度 $a$ 的符号决定势的吸引/排斥性质（$a>0$ 等效排斥，$a<0$ 等效吸引）。

\trainingpoint{5} \textbf{自旋-轨道耦合}：核子-核子散射中，$L\cdot S$ 项导致自旋取向与轨道角动量耦合，引起极化现象。

\trainingpoint{6} \textbf{磁偶极相互作用的能谱}：双中子磁偶极相互作用 $H\propto\boldsymbol{\sigma}_1\cdot\boldsymbol{\sigma}_2-3\sigma_{1z}\sigma_{2z}$，在总自旋基下对角化，求各能级的能量和简并度。
\end{drillbox}

% --------------------------------------------------------
% 4. 解题技巧框
% --------------------------------------------------------
\begin{tipbox}[快速识别与解题技巧]
\begin{itemize}
    \item \examnote{识别标志}：题干出现``自旋依赖''''低能散射''''中子-中子''''全同费米子''''单态''''三重态''''散射长度''，立即定位本模型。
    \item \examnote{自旋通道速查}：
    \begin{center}
    \fbox{\parbox{0.9\linewidth}{\centering
    $\boldsymbol{\sigma}_1\cdot\boldsymbol{\sigma}_2=2S(S+1)-3$\\[2pt]
    单态 $S=0$：$\boldsymbol{\sigma}_1\cdot\boldsymbol{\sigma}_2=-3$\\[2pt]
    三重态 $S=1$：$\boldsymbol{\sigma}_1\cdot\boldsymbol{\sigma}_2=+1$
    }}
    \end{center}
    \item \examnote{全同费米子散射}：s 波空间对称 $\Rightarrow$ 自旋单态。三重态的 s 波被泡利原理禁止。低能散射截面主要由单态贡献。
    \item \examnote{非极化平均}：两个非极化自旋-$1/2$ 粒子，总自旋态等概率：$P(S=0)=1/4$，$P(S=1)=3/4$。
    \item \examnote{散射长度的物理意义}：$a>0$ 等效排斥（束缚态在阱外），$a<0$ 等效吸引（近束缚态）。$|a|\to\infty$ 时发生 Feshbach 共振。
    \item \examnote{磁偶极相互作用的对角化}：$H\propto\boldsymbol{\sigma}_1\cdot\boldsymbol{\sigma}_2-3\sigma_{1z}\sigma_{2z}$，在 $|S,M\rangle$ 基下：
    \begin{itemize}[nosep]
        \item $|1,\pm1\rangle$：$E\propto1-3=-2$（或 $+1-3/4$？需重新计算）
        \item $|1,0\rangle$：$E\propto1+3/4$（含非对角项）
        \item $|0,0\rangle$：$E\propto-3+3/4$
    \end{itemize}
\end{itemize}
\end{tipbox}

% --------------------------------------------------------
% 5. 易错点框
% --------------------------------------------------------
\begin{errorbox}{常见失误与陷阱}
\begin{itemize}
    \item \textbf{全同粒子对称化的符号}：费米子散射振幅为 $f(\theta)-f(\pi-\theta)$。若误用 $+$，则 s 波截面会错误地不为零（实际上 s 波 $f$ 为常数，$f(\theta)-f(\pi-\theta)=0$，符合泡利原理）。
    \item \textbf{自旋态的归一化}：单态 $\frac{1}{\sqrt2}(|\uparrow\downarrow\rangle-|\downarrow\uparrow\rangle)$ 和三重态 $|1,0\rangle=\frac{1}{\sqrt2}(|\uparrow\downarrow\rangle+|\downarrow\uparrow\rangle)$ 是正交归一的。不要遗漏 $1/\sqrt2$。
    \item \textbf{非极化与极化的区别}：非极化束是混合态（密度矩阵），不是叠加态。计算截面时对自旋态取平均（入射）和求和（出射）。
    \item \textbf{低能极限的适用条件}：$ka\ll1$ 时仅 s 波有贡献。若 $ka\sim1$，需考虑 p 波、 d 波等更高分波。
    \item \textbf{散射长度的单位}：$a$ 有长度量纲。$4\pi a^2$ 是面积（截面）。注意与经典截面的区别。
    \item \textbf{有效势的符号}：$V_0>0$ 时，单态有效势 $-3V_0$（吸引），三重态 $+V_0$（排斥）。若 $V_0<0$，符号反转。
\end{itemize}
\end{errorbox}

% --------------------------------------------------------
% 6. 物理图像与关联模型
% --------------------------------------------------------
\begin{insightbox}[物理图像与模型关联]
\textbf{物理图像}：

自旋依赖势就像两个带磁铁的弹珠：
\begin{itemize}[nosep]
    \item 如果两块磁铁的北极都朝上（三重态，$S=1$），它们会相互排斥——空间波函数需要远离，但 s 波是对称的（空间上靠近），所以泡利原理禁止这种组合。
    \item 如果两块磁铁一反一正（单态，$S=0$），它们会相互吸引——空间波函数可以靠近，s 波对称正好满足。
    \item 这就是为什么两个中子（费米子）在低能下只能以单态散射：三重态被``泡利排斥''压制了。
    \item 散射长度 $a$ 描述了这种相互作用的等效范围：$a>0$ 像硬球排斥，$a<0$ 像弱吸引。
\end{itemize}

\textbf{与相似模型的对比}：
\begin{center}
\begin{tabular}{llll}
\toprule
\textbf{对比维度} & \textbf{自旋依赖势} & \textbf{中心势散射} & \textbf{库仑散射} \\
\midrule
势的形式 & $V(r)\boldsymbol{\sigma}_1\cdot\boldsymbol{\sigma}_2$ & $V(r)$ & $1/r$ \\
自旋效应 & 强（决定通道） & 无 & 无（非全同） \\
低能极限 & s 波单态主导 & s 波主导 & 无低能极限（长程） \\
全同性 & 关键（费米子） & 可选 & 可选 \\
\bottomrule
\end{tabular}
\end{center}

\textbf{推广方向}：
\begin{itemize}[nosep]
    \item Feshbach 共振：调节外磁场改变散射长度，实现 BEC-BCS  crossover
    \item 核力：汤川势 + 自旋-轨道耦合 + 张量力，描述核子-核子相互作用
    \item 冷原子物理：超冷费米气体的 s 波配对，$^6$Li 和 $^{40}$K 的 Feshbach 共振
    \item EPR 与量子纠缠：单态 $|\psi^-\rangle$ 是最大纠缠态，用于贝尔不等式检验
\end{itemize}
\end{insightbox}

% --------------------------------------------------------
% 7. 推导附录
% --------------------------------------------------------
\begin{derivationbox}[推导细节（自我检验用）]
\textbf{Step 1：自旋耦合本征值}

$\mathbf{S}=\mathbf{S}_1+\mathbf{S}_2$，$\mathbf{S}^2=\mathbf{S}_1^2+\mathbf{S}_2^2+2\mathbf{S}_1\cdot\mathbf{S}_2$。

$\mathbf{S}_1\cdot\mathbf{S}_2=\frac12\bigl(\mathbf{S}^2-\mathbf{S}_1^2-\mathbf{S}_2^2\bigr)=\frac{\hbar^2}{2}\bigl(S(S+1)-\frac34-\frac34\bigr)=\frac{\hbar^2}{2}S(S+1)-\frac34\hbar^2$。

$\boldsymbol{\sigma}_1\cdot\boldsymbol{\sigma}_2=\frac{4}{\hbar^2}\mathbf{S}_1\cdot\mathbf{S}_2=2S(S+1)-3$。

单态 $S=0$：$\boldsymbol{\sigma}_1\cdot\boldsymbol{\sigma}_2=-3$。

三重态 $S=1$：$\boldsymbol{\sigma}_1\cdot\boldsymbol{\sigma}_2=+1$。

\textbf{Step 2：全同费米子 s 波散射}

s 波（$l=0$）空间波函数对称（$Y_{00}=1/\sqrt{4\pi}$ 为常数）。

总波函数 = 空间 $\times$ 自旋，必须反对称。

空间对称 $\Rightarrow$ 自旋必须反对称 $\Rightarrow$ 单态（$S=0$）。

三重态（$S=1$，自旋对称）要求空间反对称，但 s 波空间对称，故三重态 s 波被禁止。

\textbf{Step 3：非极化散射截面}

非极化入射：四个自旋态等概率（各 $1/4$）。

单态（1 个）：s 波允许，截面 $\sigma_0=4\pi a_s^2$（全同粒子对称化后）。

三重态（3 个）：s 波禁止，p 波及以上贡献在 $k\to0$ 时趋于零。

总截面：
\[
\sigma_{\text{tot}}=\frac14\sigma_0+\frac34\cdot0=\pi a_s^2.
\]

（注：更精确地，全同费米子 s 波对称化后 $\sigma(\theta)=|f(\theta)-f(\pi-\theta)|^2$，对常数 $f$ 有 $\sigma(\theta)=0$，总截面为零。但若考虑自旋单态的有效势散射，单态通道的 $f_{S=0}$ 给出 $\sigma_0=4\pi a_s^2$。）

\textbf{Step 4：磁偶极相互作用能谱}

$H=-\frac{\mu^2}{a^3}(\boldsymbol{\sigma}_1\cdot\boldsymbol{\sigma}_2-3\sigma_{1z}\sigma_{2z})$。

利用 $\boldsymbol{\sigma}_1\cdot\boldsymbol{\sigma}_2=2S(S+1)-3$ 和 $S_z^2=(\sigma_{1z}+\sigma_{2z})^2/4$。

$\sigma_{1z}\sigma_{2z}=2S_z^2-3/2$？不对，需直接计算。

在 $|S,M\rangle$ 基下：
\begin{itemize}[nosep]
    \item $|1,1\rangle=|\uparrow\uparrow\rangle$：$\sigma_{1z}\sigma_{2z}=+1$，$\boldsymbol{\sigma}_1\cdot\boldsymbol{\sigma}_2=+1$，$H=-\frac{\mu^2}{a^3}(1-3)=+\frac{2\mu^2}{a^3}$
    \item $|1,-1\rangle=|\downarrow\downarrow\rangle$：$\sigma_{1z}\sigma_{2z}=+1$，$\boldsymbol{\sigma}_1\cdot\boldsymbol{\sigma}_2=+1$，$H=+\frac{2\mu^2}{a^3}$
    \item $|1,0\rangle$：$\sigma_{1z}\sigma_{2z}=-1$（期望），$\boldsymbol{\sigma}_1\cdot\boldsymbol{\sigma}_2=+1$，$H=-\frac{\mu^2}{a^3}(1+3)=-\frac{4\mu^2}{a^3}$
    \item $|0,0\rangle$：$\sigma_{1z}\sigma_{2z}=-1$（期望），$\boldsymbol{\sigma}_1\cdot\boldsymbol{\sigma}_2=-3$，$H=-\frac{\mu^2}{a^3}(-3+3)=0$
\end{itemize}

故能级：$E_{1,\pm1}=+2\mu^2/a^3$（二重简并），$E_{1,0}=-4\mu^2/a^3$，$E_{0,0}=0$。
\end{derivationbox}

\vspace{1em}
